\documentclass[12pt]{article}
\usepackage[utf8]{inputenc}
\usepackage[T1]{fontenc}
\usepackage{hyperref}
\title{Suku Language Resources for K-12 Curriculum Development}
\author{Compiled by Claude Sonnet 5.5 (claude-sonnet-5-5) for Andr\'as Kornai}
\date{}
\begin{document}
\maketitle

\section{Introduction}
Suku (Kisuku; ISO 639-3: sub; Glottocode: suku1259) is a Bantu language of the Democratic Republic of the Congo and Angola. Speaker counts vary by source, from about 50,000 to 80,000 or more. This report rests on a single web search and is deliberately short.

\section{Primary Educational Resources}
No primers, textbooks or school programmes for Suku were found.

\section{Linguistic and Reference Materials}
No dedicated Suku dictionary or grammar was found. The sources below are general overviews.
\begin{itemize}
\item Wikipedia, Suku language: \url{https://en.wikipedia.org/wiki/Suku_language}
\item ASJP language page (wordlist data): \url{https://asjp.clld.org/languages/SUKU_2}
\end{itemize}

\section{Religious and Community Resources}
Joshua Project reports Bible portions from 1973--1977 but no New Testament or complete Bible.
\begin{itemize}
\item Joshua Project language page: \url{https://www.joshuaproject.net/index.php/languages/sub}
\end{itemize}

\section{Digital and Audio Resources}
\begin{itemize}
\item Global Recordings Network, ``Kisuku \&amp; Kiyaka'' evangelistic and Bible-teaching audio (it may overlap with the Yaka language): \url{https://globalrecordings.net/ja/language/1543}
\end{itemize}

\section{Curriculum Development Context}
Very little educational material was found. A curriculum would have to begin with lexicon and orthography work, and with the existing Bible portions and recordings as seed texts.
\end{document}